% from main.tex
%%%%%%%%%%  scan idx 253  |  folio 240  %%%%%%%%%%
\textbf{Corollaire 2.3} (EGA III 2.2.1). Soient \(S\) un schéma quasi-compact, \(f:X\to Y\) un \(S\)-morphisme projectif de schémas de type fini sur \(S\), \(L=i^*\mathcal O_P(1)\) le faisceau inversible défini par un plongement \(i:X\to P=P(F)\), où \(F\) est un \(\mathcal O_Y\)-Module quasi-cohérent de type fini. Pour \(M\in\Ob D(X)\) et \(n\in\mathbb Z\), posons \(M(n)=M\otimes^{\mathbf L} L^{\otimes n}\). Alors, si \(E\in\Ob D^+(X)\) est pseudo-cohérent rel. à \(S\) et acyclique en degré \(>a\) \((a\in\mathbb Z)\), il existe un entier \(N\) tel que, pour tout \(n\geq N\), \(\mathbf Rf_*E(n)\) vérifie les mêmes conditions.

\emph{Preuve.} On se ramène trivialement, comme dans la preuve de (2.2 (ii)), à la situation de (2.2.2) : \(Y=S\), \(f:X=\mathbb P^r_Y\to Y\) la projection canonique, \(L=\mathcal O_X(1)\) le faisceau canonique, \(E\in\Ob D^b(X)\) un complexe pseudo-cohérent et acyclique en degré \(\geq a\), et l'on gagne grâce à (2.2.2).

\textbf{Corollaire 2.4.} Soient \(S\) un schéma quasi-compact, et \(f:X\to Y\) un \(S\)-morphisme propre de schémas de type fini sur \(S\). Si \(S\) est noethérien ou si \(f\) est projectif, le foncteur \(\mathbf Rf_*:D^+(X)\longrightarrow D^+(Y)\) induit un foncteur
\[
  \mathbf Rf_*:D^b(X)_{S\text{-}\coh}\longrightarrow D^b(Y)_{S\text{-}\coh}.
\]

\emph{Preuve.} Conséquence de (2.2) et du fait qu'il existe alors (EGA III 1.4.12 et IV 1.7.21) un entier \(N\) tel que \(R^if_*(E)=0\) pour tout \(i>N\) et tout \(\mathcal O_X\)-Module quasi-cohérent \(E\).

\textbf{Corollaire 2.5.} Soient \(f:X\to Y\) un morphisme propre et pseudo-cohérent (1.2), et \(E\in\Ob D^+(X)\). Supposons \(Y\) localement noethérien ou \(f\) projectif. Alors, si \(E\) est pseudo-cohérent (au sens absolu), il en est de même de \(\mathbf Rf_*(E)\). Si \(Y\) est quasi-compact, le foncteur \(\mathbf Rf_*\) induit un foncteur
\[
  \mathbf Rf_*:D^b(X)_{\mathrm{coh}}\longrightarrow D^b(Y)_{\mathrm{coh}}.
\]

